\section{Discussion on Non-uniformity}\label{appsec:non_uniformity}


\emph{Non-uniform} computation models allow a different program for each different input length, like boolean circuits. However, the complexity class defined by circuits can also be uniform, if we add additional assumption on the correlation between circuits of different input lengths, \emph{e.g.}, one can require the circuits for input length $n$ can be generated by a Turing Machine taken $n$ as input in using a certain amount of time and space. 

The complexity class $\Cot$ introduced in this paper can also be made uniform by enforcing an additional assumption, that the parameters of the transformer can be generalized by some Turing Machine given the input sequence length $n$. It is well-known that one can simulate the execution of the Turing Machine for any $T$ steps by a family of uniform boolean circuits of size $O(T^2)$. Thus if we enforce the parameters of transformers in $\Cot$~ to be uniform, our main theorem would imply that constant-depth transformers with uniform parameters and polynomially many steps of chain of thoughts can solve all problems in $\P$. Also note that the inference of transformers can also be done in polynomial time, we conclude it is exactly equal to $\P$.


One natural question about non-uniformity is that \emph{whether having a different transformer for each input sequence length is practical,} given that a significant portion of previous theoretical works on transformer expressiveness focuses on the uniform setting. This problem is kind of ill-defined because we haven't been able to scale up the input length to arbitrary length in practice, and thus it is not clear if it is necessary to keep scaling up the size of LLMs for longer input sequence length. But at least for the LLMs that have been seen in practice, it seems quite common to scale up the model size when dealing with longer input sequence length. Also taking the GPT architecture~\citep{radford2019language} that we focus on in this paper, having more trainable parameters is necessary for longer input sequence length, due to the trainable absolute position encoding. 

Still, one needs to note that there is a difference between natural language tasks and complexity class, where the former has a lot of memorization and does not require a strong ability to solve math problems of any sequence length. In contrast, to learn this complexity class like the composition of permutation of any length, transformers need to have the ability of \emph{length generalization}, which does seem impossible for certain non-uniform models, \emph{e.g.}, like GPT architectures with trainable absolute position encoding, because there is no way to learn the position encoding at an unseen position in the training dataset. Of course, length generalization would still be possible if GPT architecture learned the ground truth without using the trainable position encoding at all.


